\documentclass[10pt]{article}

\usepackage[letterpaper,margin=0.5in]{geometry}
\usepackage[T1]{fontenc}
\usepackage[scaled=0.95]{helvet}
\renewcommand{\familydefault}{\sfdefault}
\usepackage{microtype}
\usepackage[table]{xcolor}
\usepackage{graphicx}
\usepackage{tabularx}
\usepackage{array}
\usepackage{booktabs}
\usepackage{enumitem}
\usepackage[most]{tcolorbox}
\usepackage{setspace}

% Agency-style palette and document rhythm.
\definecolor{agencyblue}{HTML}{005A8B}
\definecolor{deepsea}{HTML}{143B5D}
\definecolor{labline}{HTML}{A9BBC8}
\definecolor{panelbg}{HTML}{F6F8FA}
\definecolor{titlebg}{HTML}{EAF2F7}
\definecolor{textgray}{HTML}{2F3B45}

\pagecolor{white}
\pagestyle{empty}
\color{textgray}
\setstretch{0.99}
\setlength{\parindent}{0pt}
\setlength{\parskip}{2.5pt}
\setlength{\tabcolsep}{4pt}
\renewcommand{\arraystretch}{1.08}
\setlist[itemize]{leftmargin=1.15em,labelsep=0.5em,itemsep=2pt,topsep=2pt,parsep=0pt}
\setlist[enumerate]{leftmargin=1.25em,labelsep=0.5em,itemsep=2pt,topsep=2pt,parsep=0pt}
\newcolumntype{Y}{>{\raggedright\arraybackslash}X}

\newtcolorbox{labbox}[2][]{
  colback=panelbg,
  colframe=labline,
  colbacktitle=titlebg,
  coltitle=deepsea,
  fonttitle=\bfseries\small,
  title=\MakeUppercase{#2},
  boxrule=0.5pt,
  arc=0pt,
  left=7pt,
  right=7pt,
  top=5pt,
  bottom=5pt,
  toptitle=3pt,
  bottomtitle=3pt,
  #1
}

\begin{document}
\small

\begin{tcolorbox}[
  colback=white,
  colframe=deepsea,
  boxrule=0.9pt,
  arc=0pt,
  left=8pt,
  right=8pt,
  top=8pt,
  bottom=7pt
]
\begin{minipage}[c]{0.17\textwidth}
\centering
\includegraphics[width=1.00\linewidth]{logo_small.png}
\end{minipage}\hfill
\begin{minipage}[c]{0.78\textwidth}
{\scriptsize\textbf{\color{agencyblue}W.E.T. LAB FIELD OPERATIONS UNIT}}\par
\vspace{2pt}
{\Large\bfseries\color{deepsea}Station 744 Dolphin Trainer Assay}\par
\vspace{2pt}
{\small Cetacean Reinforcement Trial | Behavioral Conditioning Protocol}
\end{minipage}

\vspace{7pt}
{\footnotesize
\rowcolors{1}{titlebg}{white}
\begin{tabularx}{\textwidth}{@{}>{\bfseries\color{deepsea}}l Y >{\bfseries\color{deepsea}}l Y@{}}
Program & Cetacean behavior and reinforcement testing & Status & Internal use only \\
Required subject & Assigned Dolphin role participant & Award & \$300 research funding \\
\end{tabularx}}
\end{tcolorbox}

\vspace{6pt}

\begin{minipage}[t]{0.485\textwidth}
\vspace{0pt}
\begin{labbox}{Study Objective}
Guide a dolphin subject through a sequence of target behaviors using only minimal reinforcement cues and efficient team coordination.
\end{labbox}

\begin{labbox}{Experimental System}
\begin{itemize}
\item Research lab = team
\item Dolphin = designated subject organism
\item Trainers = all remaining teammates
\item Task cards = target behavior sequence
\item Fastest successful completion = winning trial
\end{itemize}
\end{labbox}

\begin{labbox}{Protocol}
\begin{enumerate}
\item Confirm that one team member with the Dolphin role serves as the subject.
\item Draw three task cards at the station; each team may swap one card with the opposing team before the start.
\item Trainers may communicate only with approved conditioning phrases.
\item First team to complete all three task cards successfully wins the assay.
\end{enumerate}
\end{labbox}

\begin{labbox}{Operational Setup}
\begin{itemize}
\item Two labs compete head-to-head.
\item The dolphin performs while teammates act as trainers.
\item Common station props may include hats, sunglasses, a baton, a fan, and percussion items.
\item Tasks often require interaction with nearby specimens in the room.
\end{itemize}
\end{labbox}
\end{minipage}\hfill
\begin{minipage}[t]{0.485\textwidth}
\vspace{0pt}
\begin{labbox}{Conditioning Signals}
\begin{tabularx}{\linewidth}{@{}>{\bfseries}l Y@{}}
\toprule
Cue & Research equivalent \\
\midrule
Good dolphin & Positive reinforcement \\
Bad dolphin & Corrective feedback \\
Task card & Target behavior \\
Finished task & Successful conditioning event \\
\bottomrule
\end{tabularx}
\end{labbox}

\begin{labbox}{Task Deck Profile}
\begin{itemize}
\item Social behaviors: high-fives, compliments, leading a short parade.
\item Object work: baton exchange, gathering two objects, hat or sunglasses transfer.
\item Performance behaviors: dramatic entrance, freeze frame, hand-fan dance, maraca beats.
\item Communication behaviors: deliver a message, recruit a new person, copy another movement.
\end{itemize}
\end{labbox}

\begin{labbox}{Observational Notes}
Watch for rapid cue interpretation, trainer coordination, subject improvisation, and whether the dolphin recruits bystanders more efficiently than the opposing lab.
\end{labbox}

\begin{labbox}{Compliance Notes}
Only approved cue phrases may be spoken by trainers. Physical coercion is not permitted; tasks must be completed through signaling, gesture, and dolphin initiative.
\end{labbox}
\end{minipage}

\vspace{6pt}

\begin{labbox}{Outcome}
Winning lab demonstrates \textbf{superior behavioral conditioning, cue discipline, and cetacean-handler coordination}. Award disbursement: \textbf{\$300 research funding}.

\vspace{3pt}
\textbf{Completion bonus reminder:} labs that complete all mini-games earn an additional \textbf{+\$100}.
\end{labbox}

\vfill
\begin{center}
\footnotesize\color{agencyblue}\textbf{W.E.T. LAB | Station 744 | Internal Use Only}
\end{center}

\end{document}
